% BatteryModule
\section{【基本】BatteryModule — ADC + 移動平均フィルタ}

\begin{patternbox}{BatteryModule で学べるパターン}
\begin{itemize}
    \item ModuleTimerによる周期サンプリング
    \item リングバッファを使った移動平均フィルタ
    \item ADC値の電圧変換（分圧回路対応）
    \item isValidフラグによる「データが十分溜まったか」の表現
\end{itemize}
\end{patternbox}

\subsection{Config/Data構造体}

\begin{lstlisting}[style=cppstyle, caption={BatteryModule.h — Config/Data}]
static constexpr uint8_t BATTERY_FILTER_SIZE = 16;

struct BatteryConfig {
    uint8_t       adcPin;
    float         voltageDivider;      // 分圧比（例: 2.0 = 1/2分圧）
    float         adcRefVoltage;       // ADC基準電圧 [V]
    uint16_t      adcResolution;       // ADC分解能（例: ESP32は4095、AVRは1023）
    unsigned long sampleIntervalMs;
    float         lowVoltageThreshold; // 低電圧警告閾値 [V]
};

struct BatteryData {
    float voltage    = 0.0f;   // フィルタ済み電圧 [V]
    float rawVoltage = 0.0f;   // 最新の生電圧 [V]
    bool  isValid    = false;  // サンプルが十分溜まったらtrue
    bool  isLow      = false;  // 低電圧警告フラグ
};
\end{lstlisting}

\begin{pointbox}[Configにハードウェア特性を含める]
\texttt{voltageDivider}（分圧比）や\texttt{adcRefVoltage}（基準電圧）はハードウェア構成に依存する値。
これらをConfigに含めることで、異なる回路構成でも同じモジュールコードを再利用できる。
\end{pointbox}

\subsection{リングバッファと移動平均}

\begin{lstlisting}[style=cppstyle, caption={BatteryModule.cpp — updateInput()}]
void BatteryModule::updateInput(SystemData& data) {
    if (_sampleTimer.getNowTime() < _config.sampleIntervalMs) {
        return;  // 周期未到達なら何もしない（ノンブロッキング）
    }
    _sampleTimer.setTime();

    // ADC読み取り -> 電圧変換
    uint16_t adcRaw = analogRead(_config.adcPin);
    float voltage = _adcToVoltage(adcRaw);

    // リングバッファに追加
    _samples[_sampleIndex] = voltage;
    _sampleIndex = (_sampleIndex + 1) % BATTERY_FILTER_SIZE;
    if (_sampleCount < BATTERY_FILTER_SIZE) {
        _sampleCount++;
    }

    // SystemDataに書き込み
    data.battery.rawVoltage = voltage;
    data.battery.voltage    = _calcAverage();
    data.battery.isValid    = (_sampleCount >= BATTERY_FILTER_SIZE);
    data.battery.isLow      = (data.battery.voltage < _config.lowVoltageThreshold);
}
\end{lstlisting}

\begin{pointbox}[isValidの使い方]
\texttt{isValid}は「移動平均のサンプルが\texttt{BATTERY\_FILTER\_SIZE}に達したか」を示す。
起動直後はサンプルが少なく平均値が不正確なため、ロジックフェーズは\texttt{isValid == true}になるまで
電圧値を信用しないようにできる。
仕様書の「デフォルト値は正常値と区別できる無効値」ルールと同じ考え方。
\end{pointbox}

\begin{pointbox}[周期サンプリングの定石]
\texttt{updateInput()}の冒頭で\texttt{getNowTime() < interval}なら即\texttt{return}するパターンは、
周期実行の定石。ガード節として最初に書くことで、残りのコードのインデントが深くならない。
\end{pointbox}
